% =========================
% B) Benchmarks / Problem Instances
% =========================
\subsection{Benchmarks / Problem Instances}
\label{subsec:benchmarks}

\paragraph{Benchmark instances.}
We evaluate on benchmark instances defined by the tuple
\[
    \text{dataset} \times \text{split} \times \text{training recipe}.
\]
All methods share the same instances and the same evaluation pipeline:
\begin{itemize}
    \item \textbf{Datasets and splits:} \DatasetsSplits{}.
    \item \textbf{Preprocessing:} \PreprocessingDetails{}.
    \item \textbf{Training recipe (fixed globally):} optimizer, learning-rate schedule, batch size, number of epochs/steps,
    regularization, early-stopping policy (if any), and evaluation cadence are fixed across methods.
\end{itemize}

\paragraph{Objective functions.}
\textbf{Single objective.}
We define
\begin{equation}
    f(\lambda)
    = \mathcal{L}_{\mathrm{val}}\!\left(
        M_{\lambda, s_{\mathrm{eval}}}
      \right),
\end{equation}
where \(M_{\lambda, s_{\mathrm{eval}}} = \mathrm{Train}(\lambda; \mathcal{D}_{\mathrm{train}}, s_{\mathrm{eval}})\) denotes the trained model
obtained by instantiating the model specified by \(\lambda\) and optimizing its parameters on \(\mathcal{D}_{\mathrm{train}}\) under the fixed training recipe.
\(\mathcal{L}_{\mathrm{val}}(\cdot)\) evaluates this trained model on \(\mathcal{D}_{\mathrm{val}}\).

\textbf{Multiobjective (if applicable).}
When optimizing multiple objectives, we define the vector-valued objective
\begin{equation}
    \mathbf{F}(\lambda) = \big(f_1(\lambda), \ldots, f_m(\lambda)\big),
\end{equation}
where \(f_1\) is the primary objective \PrimaryObjective{} and the remaining objectives are \SecondaryObjectives{}.
When Pareto metrics require normalization (e.g., hypervolume), we apply a fixed monotone normalization policy:
\ObjectiveNormalizationPolicy{}.

\paragraph{Search space.}
All methods optimize over an identical, explicitly defined search space \(\Lambda\):
\begin{itemize}
    \item \textbf{Hyperparameters and ranges/types:} \SearchSpaceList{}.
    \item \textbf{Conditional parameters / constraints:} \SearchSpaceConditionals{}.
    \item \textbf{Canonical encoding:} \SearchSpaceEncoding{}.
\end{itemize}

\paragraph{Representativeness rationale.}
The benchmark suite is selected to represent \ClaimedDomain{} by including (where applicable):
(i) at least one small/fast instance enabling many repetitions (noise characterization),
(ii) at least one realistic instance reflecting intended deployment scale,
and (iii) heterogeneous instances to test robustness.
Conclusions are restricted to these regimes.
